% Draft paragraph for Paper 58 (m != 0 exchange); not inserted. Source: debug/pachucki_m_proof_memo.md.
\textbf{[SYMBOLIC PROOF]} The second-order equation for the weighted masters
is a theorem, given Pachucki's first-order equation in the $r_{12}$
exponent $w_1$ (his Eq.~(16)~\cite{pachucki2009}), which we use as
published.  For exponents $u_2,u_3,w_2,w_3>0$ and $|w_1|<\delta$,
$\delta<\tfrac12\min(u_2,u_3,w_2,w_3)$, the Hardy--Littlewood--Sobolev
inequality bounds the $w_1$-dependent master by $Kr$ uniformly, and in
prolate spheroidal coordinates it is smooth in $r>0$ and decays
exponentially.  Hence $g(t;w_1)=\int_0^\infty e^{-tr}f(r;w_1)\,dr$ is
holomorphic in $w_1$, its Taylor coefficients are the Laplace transforms
of the weighted masters, and Pachucki's equation, which holds for $w_1>0$,
holds on the whole disc by the identity theorem.  Its $w_1^m$ coefficient
is an identity in $t$ with no division.  Dividing it by $t^4$ makes every
term the transform of a continuous function: the logarithms enter as
$\ln(t+c)/t=\mathcal{L}[\ln c+E_1(cr)]$.  Lerch's theorem gives the
integrated identity for all $r\ge0$, and four derivatives on $r>0$ give the
equation, with the polynomials in $t$ (the terms supported at $r=0$)
removed exactly.  No derivative of a master at $r=0$ is ever evaluated, so
the logarithms the masters carry there do not enter.  The source is
elementary wherever $u_2\neq u_3$ and $w_2\neq w_3$, which is where the
division by $(u_2^2-u_3^2)(w_2^2-w_3^2)$ is allowed; the equation can be
differentiated in the exponents, and its decaying solution is unique.
